\chapter{Penutup}
\label{chap:penutup}

Bab ini menyimpulkan hasil penelitian terhadap tujuan pada Subbab~\ref{sec:tujuan} beserta temuan teknis utamanya, kemudian memberikan saran pengembangan di luar lingkup penelitian ini.

\section{Kesimpulan}
\label{sec:kesimpulan}

Penelitian ini membangun sistem monitoring detak jantung single-lead berbasis \textit{edge computing} yang menjalankan preprocessing sinyal ECG dan klasifikasi aritmia dengan model CNN \textit{hybrid} INT8 secara \textit{on-device} pada ESP32-S3, lalu mengirim hasil klasifikasi per detak melalui \textit{pipeline} MQTT dengan bufer lokal di PSRAM. Berdasarkan perancangan, implementasi, dan eksperimen pada Bab~\ref{chap:eksperimen}, kesimpulan terhadap ketiga tujuan penelitian diuraikan sebagai berikut.

Tujuan pertama tercapai. Preprocessing di ESP32-S3 terbukti setara dengan Python, yaitu seluruh tahap uji data acuan lulus, termasuk 11 R-\textit{peak} yang sama dalam toleransi 1 sampel dan keputusan akhir yang identik pada tujuh detak acuan. Beban rata-rata hanya 82~$\mu$s per sampel atau 3,0\% anggaran waktu tanpa sampel yang hilang, sedangkan puncak beban sesaat 33,4~ms diserap oleh ADC yang dibaca di ISR beserta antreannya. Segmentasi hasil detektor di perangkat memberi F1 0,6499 dibandingkan 0,6594 pada segmentasi anotasi, dan akuisisi sinyal tubuh menghasilkan rekaman yang memenuhi kriteria kualitas pada dua sesi perekaman.
% sumber: bab4 tab:uji-kesetaraan-tahap, tab:beban-komputasi, tab:segmentasi-deteksi, tab:kualitas-rekaman, tab:ketercapaian

Tujuan kedua tercapai sebagian. Model \textit{float32} pada DS2 yang berisi 22 pasien yang tidak terlihat saat pelatihan mencapai \textit{recall} 0,6996, F1 0,6594, dan AUC 0,9334. Kuantisasi INT8 menghasilkan model 22,94~KB dengan selisih \textit{recall} $+$0,0108, dan inferensinya di ESP32-S3 memakan 26,6~ms per detak dengan \textit{tensor arena} 12.948 byte, sehingga target kuantisasi, ukuran, latensi, dan memori terpenuhi tanpa ketergantungan pada \textit{cloud}. Target konsumsi daya per mode operasi belum dapat dinilai karena konsumsi daya belum diukur. Kelemahan model terletak pada kelas S dengan \textit{recall} 0,4390 pada DS2.
% sumber: bab4 tab:hasil-ds2, tab:delta-int8, tab:sumber-daya, subsec:hasil-tujuan2 butir Evaluasi DS2 (recall S 0,4390) dan butir Konsumsi Daya, tab:ketercapaian

Tujuan ketiga tercapai. Dengan sumber sinyal \textit{replay}, \textit{pipeline} MQTT mencapai PDSR 100,0\% dan 99,2\% dengan median latensi \textit{end-to-end} 939 dan 1.108~ms pada dua ulangan, sehingga memenuhi target PDSR 95\% dan latensi 2 detik. Ketika \textit{broker} dibekukan selama 24 detik, tidak ada detak yang hilang, \textit{backlog} terkuras paling lama 24,6 detik setelah koneksi pulih, dan jeda terbesar pada deret waktu di ThingsBoard hanya 911~ms karena setiap detak membawa \textit{timestamp} waktu kejadiannya.
% sumber: bab4 tab:target-pipeline, subsec:hasil-tujuan3 butir Putus Koneksi dan Pemulihan (hilang 0, 24,6 detik, jeda 911 ms), tab:ketercapaian

Selain ketercapaian tujuan, penelitian ini menghasilkan empat temuan teknis yang tidak terlihat dari rancangan awal.
Pertama, perbandingan keluaran perangkat dengan data acuan menangkap kesalahan yang berlangsung tanpa pesan apa pun. Operasi MEAN pada TensorFlow Lite Micro, yang menjalankan \textit{global average pooling}, menghasilkan probabilitas 0,336 di perangkat padahal 0,973 di PC, dan inferensi tetap selesai seolah-olah berhasil. Kesalahan itu hanya terlihat karena keluaran perangkat dibandingkan dengan acuan Python, sehingga model \textit{deploy} diganti dengan DepthwiseConv1D(32) berbobot 1/32 yang menghitung rerata yang sama dengan operasi konvolusi yang berjalan benar. Implikasinya, uji kesetaraan antara PC dan perangkat bukan formalitas, dan model yang akurat di PC tidak boleh dianggap benar di perangkat sebelum keluarannya dibandingkan.
% sumber: bab4 tab:perubahan-rancangan baris Model deploy (0,973, 0,336, DepthwiseConv1D(32) berbobot 1/32) dan subsec:hasil-tujuan1 butir Port ke Bahasa C (uji kesetaraan sebagai pengaman)

Kedua, augmentasi \textit{jitter} yang meniru \textit{error} detektor sendiri bekerja sebagai regularisasi. Residu detektor Pan-Tompkins pada DS1 berbentuk inti yang tajam dengan ekor yang berat, berbeda dari \textit{jitter} seragam pada Dias et al.~\cite{dias_2021}, sehingga augmentasi mengambil sampel dari residu terukur tersebut. Pada ablasi tiga \textit{seed}, F1 DS2 naik dari 0,5660 menjadi 0,6354 walaupun DS2 sendiri tidak diberi \textit{jitter}, sehingga manfaatnya bukan sekadar simulasi kondisi lapangan. Catatannya, DS2 ikut menilai ablasi ini, sehingga kenaikan tersebut berpotensi sedikit optimistis. Implikasinya, konstanta empiris dari penelitian lain sebaiknya diukur ulang pada detektor sendiri sebelum diadopsi.
% sumber: bab4 tab:ablasi baris bersih dan jitter (0,5660, 0,6354, tiga seed), subsec:hasil-tujuan2 butir Ketahanan terhadap Error Segmentasi (residu inti tajam dan ekor berat, berbeda dari jitter seragam Dias) dan butir Ablasi (DS2 tidak diberi jitter, regularisasi), butir Evaluasi DS2 (DS2 ikut menilai ablasi, berpotensi sedikit optimistis)

Ketiga, latensi \textit{end-to-end} dibatasi oleh sisi perangkat dan bukan oleh jaringan. Sekitar 87\% latensi, atau 78\% sampai 93\% menurut ulangan, berasal dari sisi perangkat, sedangkan \textit{broker} hanya menyumbang sekitar 1,1~ms, dan median latensi 939 dan 1.108~ms pada dua ulangan memenuhi target 2 detik. Ekor sebaran berbeda, karena p95 berayun antara 1.404 dan 3.304~ms yang diduga akibat kontensi kanal WiFi. Temuan ini nyaris terlewat, sebab pengukuran awal memberi median 1.907~ms dengan basis waktu perangkat yang beresolusi satu detik dan sempat menyimpulkan bahwa sisa latensi milik \textit{broker}. Kesimpulan itu dikoreksi setelah sumber waktu perangkat diganti dan \textit{broker} diukur langsung. Implikasinya, komponen yang dicurigai perlu diukur sebelum dipersalahkan, dan optimasi median latensi berikutnya berada di \textit{firmware}, bukan di jaringan.
% sumber: bab4 tab:dekomposisi-latensi dan subsec:hasil-tujuan3 butir Latensi End-to-End (87%, rentang 78% sampai 93% menurut ulangan, broker 1,1 ms, median awal 1.907 ms karena basis waktu satu detik, koreksi setelah gettimeofday), tab:target-pipeline (median 939 dan 1.108 ms, p95 1.404 dan 3.304 ms, target 2 detik), butir Latensi (kontensi kanal sebagai dugaan, belum diukur langsung)

Keempat, telemetri akuisisi dapat menyesatkan, sehingga kualitas sinyal tidak boleh dinilai dari keberadaan R-\textit{peak}. Detektor Pan-Tompkins memakai ambang adaptif terhadap sinyal yang ada, sehingga tetap melaporkan R-\textit{peak} ketika biopotensial sudah hilang. Rekaman 1643 melaporkan 180 R-\textit{peak} dalam 49,6 detik padahal ayunan ADC-nya hanya 55 \textit{count}, dan tiga dari enam rekaman gagal bahkan melaporkan \textit{lead-off} 0\%. Penentu yang benar karena itu adalah ayunan ADC absolut, sebab ayunan di bawah sekitar 150 \textit{count} berarti tidak ada biopotensial sama sekali, dan ayunan inilah yang dipakai sebagai gerbang kualitas pada jalur \textit{pipeline}.
% sumber: bab4 tab:kualitas-rekaman (rekaman 1643, ayunan 55 count, 49,6 detik) dan subsec:hasil-tujuan1 butir Akuisisi Sinyal Tubuh (180 R-peak dalam 49,6 detik, lead-off 0% pada tiga dari enam rekaman gagal, ambang sekitar 150 count, ayunan absolut sebagai gerbang kualitas)

Secara keseluruhan, Tujuan 1 dan Tujuan 3 tercapai, sedangkan Tujuan 2 tercapai pada seluruh target model dan kinerja di perangkat kecuali konsumsi daya yang belum diukur (Tabel~\ref{tab:ketercapaian}). Sistem ini berstatus \textit{proof of concept} di lingkungan non-klinis dan bukan alat diagnosis, dengan keterbatasan yang diuraikan pada Subbab~\ref{subsec:analisis-keterbatasan}.

\section{Saran}
\label{sec:saran}

Berdasarkan keterbatasan pada Subbab~\ref{subsec:analisis-keterbatasan}, saran pengembangan selanjutnya diuraikan pada butir berikut.
\begin{enumerate}
  \item \textbf{Pengukuran dan optimasi konsumsi daya.} Pembacaan INA219 yang sudah berjalan pada \textit{firmware} uji perangkat keras perlu diintegrasikan ke \textit{firmware} utama agar arus dapat diukur per mode operasi dengan beban \textit{replay} yang identik. Pengukuran dengan WiFi aktif perlu memakai setelan radio tanpa mode tidur seperti pada pengujian \textit{pipeline}, karena setelan itu diperkirakan menambah arus sekitar 30~mA. Bila daya menjadi kendala, \textit{kernel} ESP-NN dapat dipertimbangkan dengan syarat keluarannya diuji ulang terhadap data acuan, karena kesalahan diam-diam seperti pada temuan pertama dapat terulang.
  % sumber: bab4 subsec:hasil-tujuan2 butir Konsumsi Daya (INA219 terbaca di firmware uji perangkat keras, belum di firmware utama; sekitar 30 mA), tab:sumber-daya (kernel referensi tanpa ESP-NN), butir Port ke Bahasa C (MEAN hanya terungkap oleh data acuan)
  \item \textbf{Validasi pada subjek yang lebih beragam.} Rekaman sinyal tubuh perlu diperluas ke lebih banyak subjek dan sesi, sebaiknya disertai alat acuan detak jantung, agar variasi sinyal antarindividu terwakili. Pada sinyal tubuh tersebut, stabilisasi elektroda, kriteria R-R pada gerbang kalibrasi, dan kemampuan ambang ayunan 60~\textit{count} memisahkan elektroda lepas dari sinyal tubuh yang lemah juga perlu diuji.
  % sumber: bab4 subsec:analisis-keterbatasan butir 1, 2, 8, dan 9; subsec:hasil-tujuan3 butir Gerbang Kalibrasi dan Kualitas Sinyal (ambang 60 count, stabilisasi dan kriteria R-R tidak diuji dengan sinyal tubuh)
  \item \textbf{Pengujian ketahanan \textit{pipeline} pada skala penuh.} Skenario putus koneksi perlu diulang dengan memutus WiFi sungguhan pada \textit{ring buffer} PSRAM berkapasitas 27.000 entri dengan paket 50 detak, dengan \textit{backlog} yang jauh lebih panjang agar waktu pengurasan dan perilaku antrean penuh terukur di perangkat. Latensi juga perlu diukur pada titik akses yang tidak padat dengan lebih banyak ulangan agar p95 dapat dinilai tanpa pengaruh kontensi kanal.
  % sumber: bab4 subsec:hasil-tujuan3 butir Putus Koneksi dan Pemulihan (broker dibekukan, 43 dari 27.000 entri, konfigurasi awal) dan butir Latensi End-to-End (p95 bersih memerlukan titik akses tidak padat)
  \item \textbf{Peningkatan deteksi kelas S.} Dugaan bahwa periode refraktori 200~ms memblokir detak S yang prematur, sehingga \textit{sensitivity} detektor kelas S hanya 0,6494, perlu diuji dengan mengubah periode refraktori. \textit{Recall} S model juga turun dari 0,4473 pada DS2 menjadi 0,2494 dan 0,1503 pada data uji SVDB dan INCART, sehingga pertukaran \textit{learning rate} $10^{-3}$ terhadap $10^{-4}$ pada Tabel~\ref{tab:ablasi} dapat ditinjau ulang bila prioritas kelas berubah, dengan konsekuensi model dan data acuan pada \textit{firmware} dibuat ulang serta divalidasi ulang di perangkat.
  % sumber: bab4 tab:sensitivity-detektor (S 0,6494) dan butir Segmentasi di Perangkat (dugaan periode refraktori), subsec:hasil-tujuan2 butir Ablasi (recall S 0,4473, 0,2494, 0,1503, learning rate)
  \item \textbf{Perangkat \textit{wearable}.} Rangkaian \textit{breadboard} perlu dipindahkan ke PCB revisi kedua beserta jalur daya baterai agar perangkat dapat dipakai secara portabel, dan klasifikasi dapat diperluas dari biner menjadi penamaan jenis aritmia bila aturan penamaannya dapat dinilai pada data validasi.
  % sumber: model/docs/arsip/2026-10-08-pcb-rev2-adjustment.md (PCB revisi kedua); bab1 sec:batasan-masalah butir 1 dan 6 (biner, rangkaian baterai di luar lingkup); model/docs/README.md "Eksperimen yang tidak diadopsi" (aturan penamaan V/S tidak lolos sertifikasi VAL)
\end{enumerate}
